\section{Elicitation 2026.10.01}

\begin{tabular}{ll}
\textbf{Date:} & 2026.10.01 \\
\textbf{Consultant:} & Hector Hill / Gy'ax \\
\textbf{Elicitors:} & Zahra, Andrew, Matthias \\
\textbf{Location:} & TFS elicitation room \\
\textbf{Context:} & Second elicitation for LING 431/531 Field Methods \\
\textbf{Audio:} & \texttt{October\_01\_2026\_HH.WAV} \\
\end{tabular}

\subsection{Matthias [0:00-35:00]}

[0:43]: Found a first voices entry matching this, but I don't actually know its constituent parts. \url{https://www.firstvoices.com/gitsenimx/words/c3ac9f24-9731-4ea2-8290-d1b1b5e83dae}. There are a few entries using \textit{waiy} as `well', only found one where \textit{woh} is also translated as `well'.

\begin{exe}
    \ex \glll waiy woh \\
    waiy woh \\
    well well(?) \\
    `Let's get started'
\end{exe}

[1:33]: I don't actually know which part of my original prompt (\textit{disiwilaaks 'niin}, which was ungrammatical) Hector was translating. I don't hear \textit{ks} after \textit{siwilaa}, instead it comes after \textit{di-}. It makes no sense that \textit{di-} appear where it does if this is monoclausal, but I am at a loss.
 
\begin{exe}
    \ex \glll siwilaa di-ks-itxw \\
    learn di-ks-txw \\
    learn \Com{}-??-\Val{} \\
    `Learn with someone (?)'
\end{exe}


[2:14]: \textit{gilbilhl smex} doesn't work without context, you have to add context, as Hector puts it ``what were you doing, where did you see them, and why did you look for them''.
 
[3:00]: This way you are still ``not saying anything, there has to be some kind of accents before or after it.'' Hector wants to add \textit{gya'a-'y} after. 

To me, this discussion points at the fact that numerals cannot act as direct predicates of nouns to be interpreted as e.g. existential constructions. Instead, it seems that there has to be another predicating element for Hector to accept the sentence. It is possible that this is due to impoverished context, however, it seems that Hector was rejecting these outright as valid sentences.  

\begin{exe}
    \ex \glll gilbil smex sapa\tsb{g}ayt \tsb{g}an ?(gya'ay)\\
    gilbil smex sapa\tsb{g}ayt \tsb{g}an ?(gya'a-y)\\
    two bear amongst tree (see-\Fsg{}.\Ii{})\\
    `(I saw) two bears in the forest'
\end{exe}

[3:55]: Hector doesn't use \textit{t'ip\tsb{x}aat} as `two'.

[5:21]: Matthias' sentence, confirmed by Hector.

\begin{exe}
    \ex \label{ex:postgilbil} \glll gya'ayhl gilbilhl smex \\
    gya'a-y=hl gilbil=hl smex \\
    see-\Fsg{}.\Ii{}=\Cn{} two=\Cn{} bear \\
    `I saw two bears '
\end{exe}

[5:32]: Hector seems to drop \textit{hl} following \textit{gilbil}. When I said (\ref{ex:pregilbil}) with \textit{gilbilhl}, Hector specifically repeated it dropping \textit{=hl}.

\begin{exe}
    \ex \label{ex:pregilbil} \glll gilbil smex gya'ay \\
    gilbil smex gya'a-y \\
    two bear see-\Fsg{}.\Ii{} \\
    `I saw two bears '
\end{exe}

[5:37]: Hector says that postpredicative (\ref{ex:postgilbil}) means some form of ``excitment'', while prepredicative (\ref{ex:pregilbil}) can be said in neutral contexts. Could this intuition point to the fact that \textit{gya'a-y} is focused in (\ref{ex:postgilbil})?? Hector adds phonological stress on the predicate when restating the sentence.

[7:36]: 

\begin{exe}
    \ex \begin{xlist}
    \item \label{ex:kampre} \glll k'am gilbil smex gya'ay \\
    k'am gilbil smex gya'a-y \\
    only two bear see-\Fsg{}.\Ii{} \\
    `I saw only two bears '
    \item  \glll gya'ay k'am gilbil smex\\
    gya'a-y  k'am gilbil smex\\
    see-\Fsg{}.\Ii{} only two bear \\
    `I saw only two bears '
    \end{xlist}

\end{exe}

[8:00]: Both pre- and post-predicate \textit{k'am} are ``understandable''. Depends who you're talking to...

[9:43]: Unsure about the translation of \textit{\tsb{g}abi}. I haven't found an independent use of it yet, only seeing it in the constructions \textit{ndahl \tsb{g}abihl} or \textit{hindahl \textit{g}abihl}. [10:35] seems to show that you definitely need both of the \textit{hl}, so it would also be interesting to see where Hector systematically drops \textit{hl}.
\begin{exe}
    \ex \glll hindehl \tsb{g}abihl smex gya'an \\
    hinde=hl \tsb{g}abi=hl smex gya'a-n \\
    where=\Cn{} many=\Cn{} bear see-\Ssg{}.\Ii{} \\
    `How many bears did you see?'
\end{exe}

[11:18]: (\ref{ex:pregilbil}) is the ``direct'' answer to the question above. (\ref{ex:postgilbil}) is more ``telling you what I've \textit{seen}'. Can also answer with (\ref{ex:kampre}).

Answering the question with simply \textit{gilbil smex} is ``not finished, but we know there's two. [...] Kind of a nonchalant answer [...] if he really answered you it'd be wiht the first or second one [the full sentences]''.

[14:30]: Using \textcite[84]{davisMorphosyntaxGitksan2026} which defines \textit{=aa} as a sentence level enclitic, acting as a polar marker. Apparently, for neturel questions \textit{nee} is used alongside \textit{=aa}, but Hector seems fine with the translating the question with simply \textit{=aa}?
\begin{exe}
    \ex \glll gilbil smex gya'anaa \\
    gilbil smex gya'a-n=aa  \\
    two bear see-\Ssg{}.\Ii{}=\Q{} \\
    `Did you see two bears?'
\end{exe}

[16:20]:
\begin{exe}
    \ex \begin{xlist}
    \item \glll guphl smexhl maa'y\\
    gup=hl smex=hl maa'y  \\
    eat=\Cn{} bear=\Cn{} berries \\
    `The bear(s) ate berries'

    \item \glll yukwhl guphl smex maa'y\\
    yukw=hl gup=hl smex=hl maa'y  \\
    \Prog{}=\Cn{} eat=\Cn{} bear=\Cn{} berries \\
    `The bear(s) are eating berries'

    \item \glll yukwhl gipguphl smex maa'y\\
    yukw=hl gip$\sim$gup=hl smex=hl maa'y  \\
    \Prog{}=\Cn{} \Pl{}$\sim$eat=\Cn{} bear=\Cn{} berries \\
    `The bear(s) are eating the berries'
    \end{xlist}
\end{exe}

[17:40]: Now this is interesting. \textcite[80]{davisMorphosyntaxGitksan2026} shows that \textit{an} is the A-extraction marker, moreover, fn. 86 says that HH frequently deletes \textit{=t}, exactly what we see in the second example. 

What's fascinating is that we have what should be the head [\textit{gilbil smex}] separated from the RC [\textit{an guphl maa'y}] by the matrix verb \textit{gya'ay}. It seems that this strongly supports a movement analysis of [\textit{gilbil smex}] to a pre-predicative position.

Note that Hector first says the sentence without \textit{an}, but it feels like he corrects himself later since he is saying it while processing, so I will tentatively assume that \textit{an} is still necessary in this construction.

Hello it's me from future: Apparently this isn't crazy since it is still ambiguous; both focus fronting and a headless RC would show the same morphology. 

\begin{exe}
    \ex \glll ba\tsb{g}adil/gilbil smex gya'ay an guphl maa'y \\
    ba\tsb{g}adil/gilbil smex gya'a-y an gup=hl maa'y \\
    \textsc{hum}.two/two bear see-\Ssg{}.\Ii{} \Ax{} eat=\Cn{} berries \\
    `I seen two bears eating berries'
\end{exe}

[20:22]: ``It's almost out of place, but it's understandable that you've seen two bears eating berries''. If using \textit{gipguphl} it means a lot of berries.

\begin{exe}
    \ex \glll ?? gilbil smex an guphl maa'y gya'ay \\
    {} gilbil smex an gup=hl maa'y gya'a-y \\
    {} two bear \Ax{} eat=\Cn{} berries see-\Ssg{}.\Ii{} \\
    Intended: `I saw two bears eating berries'    
\end{exe}

[20:49]: In response to above, Hector offers this. 

\begin{exe}
    \ex \glll gya'ay gilbil smex an guphl maa'y \\
    gya'a-y gilbil smex an gup=hl maa'y \\
    see-\Ssg{}.\Ii{} two bear \Ax{} eat=\Cn{} berries \\
    `I saw two bears eating berries'    
\end{exe}

[21:16]: I think I hear the attributive \textit{-m} after \textit{s\tsb{g}an}.

\begin{exe}
    \ex \glll s\tsb{g}anim maa'y \\
    s\tsb{g}an-m maa'y \\
    bush-\Attr{} berries \\
    `bush/branch of berries'    
\end{exe}

[21:35]:  I swear I don't hear [=hl] here lol.

\begin{exe}
    \ex \glll miiluxw[=hl] ba\tsb{g}adil hanaa\tsb{k}' \\
    miiluxw[=hl] ba\tsb{g}adil hanaa\tsb{k}' \\
    dance[=\Cn{}] two.\Hum{} woman.\Pl{} \\
    `Two women danced'
\end{exe}

    
[23:52] + [24:59]: 

\begin{exe}
    \ex \glll gilbilhl/ba\tsb{g}adil hanaa\tsb{k}' miiluxwit  \\
    gilbil=hl/ba\tsb{g}adil hanaa\tsb{k}' miiluxw-t \\
    two.\Hum{}/two=\Cn{} woman.\Pl{} dance-\Tsg{}.\Ii{}  \\
    `Two women danced'
\end{exe}

[24:09]: 

\begin{exe}
    \ex \glll yukwhl miiluxw=hl ba\tsb{g}adil hanaa\tsb{k}' \\
    yukw=hl miiluxw=hl ba\tsb{g}adil hanaa\tsb{k}' \\
    \Prog{}=\Cn{} dance=\Cn{} two.\Hum{} woman.\Pl{} \\
    `Two women are dancing'
\end{exe}

[24:44] + [25:25]:: Both readings don't seem to work. 

\begin{exe}
    \ex \glll * ba\tsb{g}adil miiluxwhl hanaa\tsb{k}' \\
    {} ba\tsb{g}adil miiluxw=hl hanaa\tsb{k}' \\
    {} two.\Hum{} woman.\Pl{} dance=\Cn{}\\
    Intended: `Two women danced' / `Women danced twice'
\end{exe}

[25:30]: Unsure if \textit{he'ksimux} `again' can be decomposed! Comment: ``They stop/sit down and they begin again.''

\begin{exe}
    \ex \glll he'ksimux miiluxwhl ba\tsb{g}adil hanaa\tsb{k}' \\
    he'ksimux miiluxw=hl ba\tsb{g}adil hanaa\tsb{k}' \\
    again dance=\Cn{} two.\Hum{} woman.\Pl{} \\
    `Two women danced again'
\end{exe}

[26:25]:

\begin{exe}
    \ex \glll hilthl smex sapa\tsb{g}ayt \tsb{g}an \\
    hilt=hl smex sapa\tsb{g}ayt \tsb{g}an \\
    many=\Cn{} bear amongst tree\\
    `There are many bears in the forest.'
\end{exe}

[26:53]: Hector didn't confirm this, but didn't reject it either? In response Hector mentioned \textit{hiltdaya} (or something analagous) may be used as a response? To be like ``oh, there are many?''. 

For Michael: if \textit{=aa} combines with \textit{hilda}, do we get palatilization as in \textit{hildaya}?

\begin{exe}
    \ex \glll gya'a'yhl hilda smex \\
    hilt=hl smex sapa\tsb{g}ayt \tsb{g}an \\
    many=\Cn{} bear amongst tree\\
    `There are many bears in the forest.'
\end{exe}

[27:38]: Apparently there are other words for `many' that Hector hears when he goes home, but he doesn't use them. 

\begin{exe}
    \ex \glll hilthl smex gya'a'y \\
    hilt=hl smex gya'a-'y \\
    many=\Cn{} bear see-\Fsg{}.\Ii{}\\
    `I saw many bears.'
\end{exe}

[28:40]: Don't really hear a \textit{=hl} after \textit{silkw}. Is it there underlyingly?

\begin{exe}
    \ex \glll silkw \tsb{g}abiihl smex gya'a'y \\
    silkw \tsb{g}abii=hl smex gya'a-'y\\
    amazing \Cnt{}=\Cn{} bear see-\Fsg{}.\Ii{}\\
    `I saw a lot [=an amazing amount] of bears.'
\end{exe}

[29:40]: ``There has to be something else to complete that [=the utterance]''. I'm guessing it's ungrammatical? I have a suspicion that \textit{\tsb{g}abiihl} is doing some funky stuff, since otherwise weak quantifiers are attested post-predicatively?

\begin{exe}
    \ex \glll * gya'a'y silkw \tsb{g}abiihl smex \\
    {} gya'a-'y silkw \tsb{g}abii=hl smex\\
    {} see-\Fsg{}.\Ii{} amazing= \Cnt{}=\Cn{} bear \\
    `I saw a lot [=an amazing amount] of bears.'
\end{exe}


[30:48]: Not sure if this is a full sentence

\begin{exe}
    \ex \glll silkw wildit \\
    silkw wild-t \\
    amazing do-\Tsg{}.\Ii{} \\
    `He/she did something bad (unsure)'
\end{exe}

[30:50]: I don't really know the status of a. below and why Hector first said it with \textit{\tsb{g}abihl}... Totally possible that it has a use case here, but idk what it is.

\begin{exe}
    \ex 
    \begin{xlist}
        \item \glll gya'a'y \tsb{g}abiihl win silkw wilhl smex \\
    gya'a-'y \tsb{g}abii=hl win silkw wilt=hl smex \\
    see-\Fsg{}.\Ii{} \Cnt{}=\Cn{} \Comp{} amazing do-\Cn{} bear \\
    `Possibly: I saw how bad the bears act.'

    \item \glll gya'a'y win silkw wilhl smex \\
    gya'a-'y win silkw wilt=hl smex \\
    see-\Fsg{}.\Ii{} \Comp{} amazing do-\Cn{} bear \\
    `I saw how bad the bears act.'
    \end{xlist}
\end{exe}

[32:25]: Again, in response we can say something like ``hilhiltaya'' was is a lot?

\begin{exe}
    \ex \glll hilthl smex gya'a'y an guphl maa'y \\
    hilt=hl smex gya'a-'y an[-t] gup=hl maa'y \\
    many=\Cn{} bear see-\Fsg{}.\Ii{} \Sx{}[-\Tsg{}.\Ii{}] eat=Cn{} berries \\
    `I saw many bears eating berries.'
\end{exe}

[33:55]: ``It's broken, got to put gya'a'y in the beginning.''

\begin{exe}
    \ex \glll * hilthl smex an guphl maa'y  gya'a'y\\
    {} hilt=hl smex an[-t] gup=hl maa'y  gya'a-'y\\
    {} many=\Cn{} bear  \Sx{}[-\Tsg{}.\Ii{}] eat=Cn{} berries see-\Fsg{}.\Ii{}\\
    Intended: `I saw many bears eating berries.'
\end{exe}

[34:29]: Not 100\% on this because it was a bit broken up, so not the best-elicited sentence by me. But likely better than the above.

\begin{exe}
    \ex \glll gya'a'y hilthl smex an guphl maa'y \\
    gya'a-'y hilt=hl smex an[-t] gup=hl maa'y \\
    see-\Fsg{}.\Ii{} many=\Cn{} bear  \Sx{}[-\Tsg{}.\Ii{}] eat=Cn{} berries\\
    Intended: `I saw many bears eating berries.'
\end{exe}




\subsection{Zahra [35:00-59:10]}

[35:40]:\textit{}

\begin{exe}
    \ex \glll 
    hilthl smex[hl?] gya'a'y \\
    hilt=hl smex[=hl?] gya'a-'y \\
    many=\Cn bear[=\Cn?] see-\Fsg.\Ii \\
    `there are a lot of bears' (repeated from earlier)
\end{exe}

[]: hilthl t'uuts'xw-m smex  (g’a’ay) loot \\
Lots of black bears there.

[36:04]
\begin{exe}
    \ex \glll 
    hilthl smex t'uuts'xwit (wilt loot) \\
    hilt=hl smex t'uuts'xw=it (wil-t loo-t) \\
    many\Cn{} bear black=\Sx{} (\Comp-3 \Obl-3) \\
    \glt 'there are many black bears' 
\end{exe}

[36:24]
\begin{exe}
\ex \glll 
    hilt t'uuts'xw-m smex (loot) \\
    hilt t'uuts'xw-m smex (loo-t) \\
    many[=\Cn{}?] black-\Attr{} bear (\Comp-3 \Obl-3) \\
    \glt 'there are many black bears' 
\end{exe}

[36:30]
\begin{exe}
\ex \glll 
    hilt t'uuts'xw-m smex gy'a'ay loot \\
    hilt t'uuts'xw-m smex gy'a'a-y loo-t \\
    many[\Cn{}?] black-\Attr{} bear see-\Fsg.\Ii{} \Obl-3 \\
    \glt 'I see many black bears' 
\end{exe}

[37:21]: 
\begin{exe}
    \ex \glll
    (yukwhl) wokhl 'wii t'uuts'xwim smex  \\
    (yukh=hl) wok=hl 'wii t'uuts'xw-m smex  \\
    (\Prog=\Cn) sleep=\Cn{} big black-\Attr{} bear \\
    \glt `The big black bear is sleeping.'
\end{exe}

[38:20]: 
\begin{exe}
    \ex \glll
    wokhl t'uuts'xwim smex k'yoots \\
    wok=hl  t'uuts'xw-m smex k'yoots  \\
    (\Prog=\Cn) sleep=\Cn{} big black-\Attr{} bear \\
    \glt `The black bear slept yesterday.'
\end{exe}

% [38:40]:


% []: (ga’a’y) /  (yukw) hadiks=hl ['wii t'uuts'xw-it smex] (k’yoots) \\
% swim=CN big black-SX bear \\
% The big black bear is swimming.

[38:58]:
% \begin{exe}
%     \ex \glll
%     haadikshl 'wii t'uuts'xwim smex \\
%     haadiks=hl 'wii t'uuts'xw-m smex \\
%     \\
% \end{exe}

% [39:04]:
\begin{exe}
    \ex \glll
    (yukwhl) haadikshl t'uuts'xwim smex \\
    (yukw=hl) haadiks=hl t'uuts'xw-m smex \\
    (\Prog=\Cn) swim=\Cn{} black-\Attr{} bear \\
    \glt `The black bear is swimming.'
\end{exe}

[39:11]:
\begin{exe}
    \ex \glll
    gya'a'y win yukwhl haadikshl t'uuts'xwim smex  \\
    gya'a-'y win yukw=hl haadiks=hl t'uuts'xw-m smex  \\
    see-\Fsg.\Ii{} \Comp{} \Prog=\Cn{} swim=\Cn{} black-\Attr{} bear \\
\end{exe}

[39:27]:
\begin{exe}
    \ex \glll
    gya'a'y win haadikshl t'uuts'xwim smex ky'oots \\
    gya'a-'y win haadiks=hl t'uuts'xw-m smex ky'oots \\
    see-\Fsg.\Ii{} \Comp{} swim=\Cn{} black-\Attr{} bear yesterday  \\
\end{exe}


[40:00]: 
\begin{exe}
    \ex \glll
    t'uuts'xwim smex hadiksit \\
    t'uuts'xw-m smex hadiks-it \\
    black-\Attr{} bear swim-\Sx \\
    \glt `Only the black bear is swimming.'
    Comment: “If the other brown bear or the black bear is still standing; only the black bear is swimming. Even though there's other black bears there, there's only one that's swimming.”
\end{exe}

[40:36]: 
\begin{exe}
    \ex \glll 
    t'uuts'xwim smex wok-it \\
    t'uuts'xw=m smex wok-it \\
    black=\Attr{} bear sleep-\Sx{} \\
    \glt `(Only) the black bear is sleeping (orig tr.; see BDM 2017).'
\end{exe}


[41:40]: 

\begin{exe}
    \ex \glll
    gya'a'y hilda hun \\
    gya'a'y[=hl?] hild-a hun \\
    see-\Fsg.\Ii[=\Cn?] many-\Attr{} fish \\
    \glt `I saw many fish.' 
    Comment: "Answers a question ... somewhere where there’s no fish, no water. ... You’re telling a story... Might be surprising." [Later says, there of course has to be water.]
\end{exe}

[43:05]: 
\begin{exe}
    \ex \glll 
    hilthl hun gya'a'y  \\ 
    hilt=hl hun gya'a-'y  \\ 
    many=\Cn{} fish see-\Fsg.\Ii  \\
    \glt `I saw a lot of fish.'
\end{exe}

[43:38]:
\begin{exe}
    \ex \glll
    hilt nagya'ahl hun \\ 
    hilt na=gya'a=hl hun \\ 
    many-\Attr{} \Fsg.\I=see=\Cn fish \\
    \glt `I saw fish many times.'
\end{exe}

[43:48]:
\begin{exe}
    \ex \glll
    agwili gya'a'yhl hun \\ 
    agwili gya'a'y=hl hun \\ 
    always \Fsg.\I=see=\Cn fish \\
    \glt `I always see fish.'
\end{exe}

[44:25]:
\begin{exe}
    \ex \glll
    (gya'a'y) hilda hun \\
    (gya'a'y) hild-a hun \\
    see-\Fsg.\Ii[=\Cn?] many-\Attr{} fish \\
    \glt `(I saw) many fish.' 
    Comment: "Where did you see them?"
\end{exe}

[44:30]:
\begin{exe}
    \ex \glll
    g̱a'ande magya'at \\
    g̱a'a-nde ma=gya'a-t \\
    \Loc-\Wh{} \Ssg.\I=see-3.\Ii  \\
    \glt `Where did you see [them]?'
    % Comment: 
\end{exe}

[44:40]:
\begin{exe}
    \ex \glll
    gya'a'y hilda hun duu'w sus \\
    gya'a'y hild-a hun  duu'w sus[t?] \\
    see-\Fsg.\Ii[=\Cn?] many-\Attr{} fish over.there \Dem.\Dist \\
    \glt `I saw many fish over there.' 
    Comment: "Now you're pointing."
\end{exe}

\textit{sus(t)} for 'that'/'over there' in FV:Gitsenimx and glossed as \Dem.\Dist in a presentation (Matthewson, Schwan and Todorović 2023). More commonly seems to be with \textit{g̱oo(')}/\textit{g̱a'a} \Loc, which Hector uses in the next sentence in \textit{g̱a'asun} (proximal) and later \textit{g̱a'asus} (distal). See \textit{g̱oosun} from the grammar sketch. Storybook says \textit{g̱oosust} is \textit{qo:[-t]=t=xʷist} `\Loc[3.\Ii]=\Pn=\Dem.\Dist'. I'm not totally clear why\Pn{} would turn to \textit{s=} unless \textit{duu'w} has [\textit{-t}] for some reason.

[45:08]:
\begin{exe}
    \ex \glll
    gya'a'y hildihl hun g̱a'asun \\
    gya'a'y hild=hl hun  g̱a'a[-t]=t=xʷin \\
    see-\Fsg.\I many=\Cn{} fish \Loc[3.\Ii]=\Pn=\Dem.\Prox \\
    \glt `I saw many fish over there.' 
    Comment: "We all saw it, but you want to specifically ."
\end{exe}

[45:24]:
\begin{exe}
    \ex \glll
    gya'a'y hild-a hun g̱a'asun \\
    gya'a'y hild-a hun  g̱a'a[-t]=t=xʷin \\
    see-\Fsg.\I many-\Attr{} fish \Loc[3.\Ii]=\Pn=\Dem.\Prox \\
    \glt `I saw many fish over [there/here].' 
    Comment: If we're all there and all saw them, it's like you're bragging, taking credit. 
\end{exe}

[47:14]:
\begin{exe}
    \ex \glll
    gya'a'y hild-a hun g̱a'asus \\
    gya'a'y hild-a hun  g̱a'a[-t]=t=xʷis \\
    see-\Fsg.\I many-\Attr{} fish \Loc[3.\Ii]=\Pn=\Dem.\Dist \\
    \glt `I saw many fish over there.' 
    Comment: If we're all there and all saw them, it's like you're bragging, taking credit. 
\end{exe}

[47:26]:
\begin{exe}
    \ex \glll
    gya'a'y hilda hun duu'w \\
    gya'a'y hild-a hun  duu'w \\
    see-\Fsg.\Ii[=\Cn?] many-\Attr{} fish over.there \\
    \glt `I saw many fish (way) over there.' 
    Comment: "That means you were over there; you're not seeing it now; or you're over here and you still see it"
\end{exe}

[48:25]:
\begin{exe}
    \ex \glll
    amhl g̱abihl hun [g̱a'a?]susaa?  \\
    am=hl g̱abi=hl hun [g̱a'a[-t]?]=t=xʷis=aa ?  \\
    good=\Cn{} \Cnt=\Cn{} fish [\Loc[3.\Ii]?]=\Pn=\Dem.\Dist=\Q \\
    \glt `Is there enough fish over there?'
    Comment: "[My mom] was talking about fish over here, but I thought she was talking about over there." 
\end{exe}

[49:50]:
\begin{exe}
    \ex \glll
    (si-)hildihl (hun) k'ii'namin loot? \\
    (si-)hilt=hl (hun) k'ii'nam-n loo-t? \\
    (\Caus1-)many=\Cn (fish) give-\Ssg.\Ii \Obl-3.\Ii \\
    \glt `How much (fish) should I give you?'
\end{exe}

% May want to double-check \textit{si}.

[50:02]:
\begin{exe}
    \ex \textit{Context: You're putting fish in his bucket.} \glll
    hildihl (hun)? \\
    hilt=hl (hun) \\
    many=\Cn (fish) \\
    \glt `How much (fish)?'
\end{exe}

[51:08]:
\begin{exe}
    \ex \glll
    dag̱abiit? \\ 
    (da,nda?)-g̱abi-t? \\
    (\Pl,\Wh?)-\Cnt-3.\Ii? \\
    \glt `How many?'
\end{exe}

[51:15]:
\begin{exe}
    \ex \textit{Context: You're putting fish in his bucket.} 
    \glll
    hildihl (hun)! \\
    hilt=hl (hun) \\
    many=\Cn (fish) \\
    \glt `more! / lots!' (orig. tr.)
\end{exe}

[52:10]:
\begin{exe}
    \ex \glll
    sihildihl \\
    si-hilt-hl \\
    \Caus1-many\Cn \\
    \glt `make it more'
\end{exe}

[52:45]:
\begin{exe}
    \ex \glll
    sihilthl maa'y k'iinamin \\
    si-hilt-hl maa'y k'iinam-in \\
    \Caus1-many\Cn berry give-\Ssg.\Ii \\
    \glt `make sure you give a lot of berries'
\end{exe}

[52:55]:
\begin{exe}
    \ex \glll
    silkw g̱abihl smex \\
    silkw g̱abi=hl smex \\
    amazing \Cnt=\Cn{} bear \\
    \glt "there's a lot of bears"
\end{exe}

[53:05]:
\begin{exe}
    \ex \glll
    sihilthl hun k'iinamin \\
    si-hilt-hl hun k'iinam-in \\
    \Caus1-many\Cn fish give-\Ssg.\Ii \\
    \glt `make sure you give a lot of fish'
\end{exe}

[53:50]: \begin{exe}
    \ex \glll
    jehlhl Gitanyow Kispiox  \\ 
    jehl=hl Gitanyow Kispiox  \\ 
    beat=hl Gitanyow Kispiox \\
    \glt 'Gitanyow beat Kispiox; Kispiox lost to Gitanyow.' 
\end{exe}

Seems like \textit{jehl} can be (in)transitive, potentially a bit interesting.

[54:14]: \begin{exe}
    \ex \glll
    yukwhl hlithl Kispiox g̱an Gitanyow 'ii jehl Kispiox. \\ 
    yukw=hl hlit=hl Kispiox g̱an Gitanyow 'ii jehl[=hl?] Kispiox. \\ 
    \Prog=\Cn{} soccer=\Cn{} Kispiox and Gitanyow and lose[=\Cn?] Kispiox \\
    \glt 'Kispiox and Gitanyow were playing soccer, and Kispiox lost.'.
\end{exe}

\textit{hlit} 'soccer'. \textit{xsdaa} 'win'.

[54:47]
\begin{exe}
    \ex \glll
    k'i'y win jehl Kispiox ahl Gitanyow. \\ 
    k'i'y win jehl Kispiox a[-t]=hl Gitanyow. \\ 
    one \Comp{} lose Kispiox \Prep[-3.\Ii]=\Cn{} Gitanyow \\
    \glt 'Kispiox lost by one from Gitanyow'.
\end{exe}

% \begin{exe}
%     \ex \gll
%     gwile'l win jehl Kispiox[hl ?] Gitanyow. \\ 
%     one \Comp{} lose Kispiox {} Gitanyow \\
%     \glt 'Kispiox lost by one from Gitanyow'.
% \end{exe}

[55:36]:
\begin{exe}
    \ex \glll
    gwile'l win xsdaahl Gitanyow 'ii ḵ'am gilbil jebhl Kispiox. \\ 
    gwile'l win xsdaa=hl Gitanyow 'ii ḵ'am gilbil jep=hl Kispiox. \\ 
    three \Comp{} win[=score] Gitanyow and only two make=\Cn Kispiox \\
    \glt 'Gitanyow scored three times and Kispiox only scored twice'.
\end{exe}

[57:29]: 
\begin{exe}
    \ex \glll
    gwile'l win ts'ilim hlo'oxsdiithl hlit ehl Kispiox \\
    gwile'l win ts'ilim hlo'oxs-diit=hl \\
    three \Comp{} in kick-\Tpl.\Ii=\Cn{} ball \Prep[-3.\Ii]=\Cn{} Kispiox \\
    \glt `They kicked three balls into Kispiox.'
\end{exe}

Seems lilke \textit{ts'ilim} and \textit{ts'imil} are related forms with metathesis? See grammar sketch vs. FV.

[57:55]:
\begin{exe}
\ex \glll
gwile'l win xsdaahl Gitanyow ii ḵ'am gilbil jebhl Kispiox \\
gwile'l win xsdaa=hl Gitanyow ii ḵ'am gilbil jep=hl Kispiox \\
three \Comp{} win=\Cn{} Gitanyow and only two make=\Cn{} Kispiox \\
\glt `Gitanyow scored three times and Kispiox only made two.'
\end{exe}

[58:13]: 
\begin{exe}
\ex \glll
gwile'l ts'ilim hlo'oxst Gitanyaaw ii ḵ'am gilbil jebhl Kispiox \\
gwile'l ts'ilim hlo'oxs-t Gitanyow ii ḵ'am gilbil jep=hl Kispiox \\
three in kick-3.\Ii{} Gitanyow and only two make=\Cn{} Kispiox \\
\glt `Gitanyow kicked it in three times and Kispiox only made two.'
\end{exe}

I'm missing some of the ungrammatical sentences here; potentially will need to come back and transcribe those.

\subsection{Andrew [59:10-1:30:36]}